\documentclass[a4paper,twoside]{article}

\usepackage{makeidx} 
\usepackage{epsfig}

\newcommand{\method}[4]{
\noindent
 \index{#1}
 \begin{tabular}{ll}
   {\bf {\small Method:}} & {\large {\tt #1}} \\    
   {\tt {\small Arguments:}} & {\rm #2} \\ 
   {\tt {\small Result:}}    & {\rm #3}
 \end{tabular}
 \\*
 \begin{tabular}{ll}
   \hspace{1mm}&
   \begin{minipage}[l]{0.8\textwidth}
     {\sl #4} 
   \end{minipage}
 \end{tabular}
\vspace{0.5cm}
}

\title {WireAPI for Ruby}
\author {Jurgen Van Ham}

\makeindex
\begin{document}

\maketitle
\begin{figure}[htp]
\centering
\includegraphics[width=1.5cm,height=1.5cm]{ruby-logo.eps}
\hspace{1cm}
\includegraphics[width=3cm,height=1.5cm]{openamq-logo.eps}
\end{figure}

\tableofcontents



\section{What?}

OpenAMQ offers the WireAPI library to interact with 
the OpenAMQ server, since the binding to Ruby written by 
Benjamin April $<$bapril@gmail.com$>$ dating from 2007 was very 
limited, this version is mainly a fresh start to expose the 
WireAPI-library to Ruby in a more complete way and closer to the 
naming used in Ruby. Apart from names this binding tries to be
close to the C api to make it easier for users of the bindings
of other languages to use this binding also. In the rest of this 
text there is a  distinction between the 
WireAPI-library\index{WireAPI-library} and the 
WireAPI-binding\index{WireAPI-binding}. The binding is 
specific for ruby while the library belongs to the OpenAMQ project.

Around the same time as this project was released to the public
another implementation appeared. Chris Wong seems to be closer
to the ruby way of doing things. His implementation can be found
at http://github.com/lightwave/ropenamq/. Until now these bindings
are not really compared as far as I know.

The protocol itself might be easy but to understand how it works
and can be used is a matter of reading on the net.\label{sec:info-url}
\begin{itemize}
\item http://www.openamq.org/doc:amqp-background  
\item http://www.openamq.org/doc:user-1-introduction  
\item http://rajith.2rlabs.com/category/amqp/ - details about using 
  and concepts are explained very well here.
\item http://www.openamq.org/tutorial:content-based-routing - how to use
  routing with headers
\end{itemize}


Wireshark is aware of the AMQP protocol so it can be inspected in a
similar way as http.



\section{Alternatives}

The WireAPI library is a part of OpenAMQ but there are other
servers. Since the protocol is open this must be also possible
to use the WireAPI with those other servers. The AMQP
website\footnote{http://jira.amqp.org/confluence/display/AMQP/AMQP+Products} 
metions the following implemetations

\begin{itemize}
  \item OpenAMQ -  http://www.openamq.org/
  \item Apache QPid - http://cwiki.apache.org/qpid/
  \item RabbitMQ - http://www.rabbitmq.com/
  \item $\oslash$MQ - http://www.zeromq.org/
\end{itemize}

Apache-QPid\index{QPid} offers a Ruby client that is written 
only in Ruby without any native code. There are 2 versions of 
the server for this project. A C++ and a Java implementation which
seem to have different objectives.

The WireAPI library claims to be the only one that supports direct 
mode for higher performance at this moment.

RabbitMQ is written in Erlang which is nice for multiprocessor 
scalability. But they don't seem to offer a Ruby client.

ZeroMQ and OpenAMQ are both created by IMatix, ZeroMQ can support 
AMQP wire protocol, so probably this binding will work with it.



\section{Build the GEM}

For the development of this binding we used OpenAMQ-1.3, it 
doesn't mean at all other versions won't work. Just this 
version worked during the testing.

Make sure the {\tt OpenAMQ-1.3c5.tar.gz} with a an environment
variable {\tt IBASE} for example /opt/amq is installed. If this 
is not yet available, install it, while being very patient. The 
variable {\tt AMQ\_HOME} needs to have the value used as IBASE 
to be installed. The script {\tt build\_gem.sh} will create a 
gem file that can be installed similar to other gems. 

During the {\tt gem install} this AMQ\_HOME variable is needed to
find the libraries and include files. A test for this variable is 
in the ext\_conf.rb file.

On some systems there is a problem with the definition of the
type 'off64\_t' in the file apr.h, if this happens replace it
with '\_\_offset64\_t' so put a double underscore as a prefix
which can help.

This manual is written in \LaTeX { }  and can be build with the
included makefile file {\tt Makefile.doc}. The 'normal' makefile
is generated with by the {\tt ext\_conf.rb} script and applies
only to the code for the binding.

For instance:
\begin{verbatim}
make -f Makefile.doc manual.pdf
\end{verbatim}


\section{Structure of the binding}

The WireAPI is a class\footnote{There can be multiple instances
but they will point to the same object.} but mainly acts as a 
namespace, there are a few classes in this namespace. 

\begin{itemize}
\item WireAPI::AMQ\_Auth used only to create a connection
\item WireAPI::AMQ\_Connection - In theory a connection could
  contain multiple sessions, but OpenAMQ doesn't support this
  anymore.
\item WireAPI::AMQ\_ContentBasic - The content or message that
  will be used to interact via the server.
\item WireAPI::AMQ\_Error - used for exceptions from the 
  WireAPI-library or the binding
\item WireAPI::AMQ\_ServerError - used for errors that were sent
  by the server
\item WireAPI::AMQ\_Session - In the future session and connection 
  will be merged together, so many methods  are available in both now.

\item WireAPI::AMQ\_FieldList - contains headers of a AMQ\_ContentBasic. 
\item WireAPI::AMQ\_int64 - used for delivery tags, for ruby this
  object is opaque since there is no real need to see the content
  of it at this moment.
\end{itemize}


Most objects have a {\tt destroy} method to force the destruction 
of the internal object in the WireAPI binding. The garbage collector 
will call this {\tt destroy} method when needed. After calling 
{\tt destroy} on a ruby object of the binding it should no longer
be used. The Ruby objects are a wrapper around the objects in the
library, when their internal core is removed these objects are have
no meaning anymore.


\section{Concepts}

The aim of this section is not to explain all details about 
AMQP but to give a bit of quickstart info enought to understand
what the main concepts are. Since the words might be similar
to other system like for instance JMS \index{JMS} the sematics
are very different. For example \emph{topic} and \emph{queue} 
are not related at all to the JMS concepts, so try not to 
find any relationship with other system to understand AMQP 
in the first place. The websites mentioned in section 
\ref{sec:info-url} (page \pageref{sec:info-url}), and the 
website of AMQP will offer more information.

\subsection{Durable}

The word has two meanings, the first is that queues and exchanges
can be durable in AMQP (but not in OpenAMQ), it means the server
will create them again after a restart.

A \emph{durable subscription} \index{durable} is the opposite of a 
\emph{transient subscription} \index{transient}it means that 
after a consumer stops the received data stays in the queue. To 
make a transient subscription use the exclusive flag (which implies 
the  auto\_delete flag) when on declaring the queue in the 
consumer program. When the consumer stops the queue will 
be deleted as its content.



\subsection{Queue and Exchange}\index{Queue}\index{Exchange}
An Exchange is an inbound message router, and a 
Queue is a message delivery vehicle. Messages come 
into an Exchange, which directs them to Queues. Each 
Exchange is responsible for routing messages to the 
consumers that want them, and uses the logic embedded 
in it to support this.

\subsection{Binding}\index{Binding}

A binding defines the relationship between an 
exchange and a queue. In other words it defines 
the routing criteria. The most simple case is where 
the binding equals the queue name. A binding 
decouples a queue from an exchange. The same 
queue can be bound to any number of exchanges 
using the same criteria or different criteria. 
Different queues can be bound to the same exchange 
using the same routing criteria as well.



\subsection{Exchange types}

The type of an exchange defines how the routing
will work for instance by routing\_key, headers
or no control at all. OpenAMQ has already some 
predefined exchanges so in most cases declaring 
your own exchange might not be
needed. 

\begin{itemize}
  \item Direct Exchange \index{direct} - 
    The exchange does a direct  match between 
    the routing key provided in the 
    message and the routing criteria used when a 
    queue is bound to this exchange.    
  \item Topic Exchange \index{topic} - 
    The exchange does a wildcard 
    match between the routing key and the routing 
    pattern specified in the binding. The routing 
    key is treated as zero or more more words, 
    delimited by ‘.’ and supports special wildcard 
    characters. “*” matches a single word and 
    ‘\#’ matches zero or more words.
  \item Fanout Exchange \index{fanout} - 
    Queues are bound to this 
    exchange with no arguments. Hence any message 
    sent to this exchange will be forwarded to all 
    queues bound to this exchange.
  \item Headers Exchange \index{headers} - 
    Queues are bound to this 
    exchange with a table of arguments containing 
    headers and values (optional). A special argument 
    named “x-match” determines the matching algorithm, 
    where “all” implies an AND (all pairs must match) 
    and “any” implies OR (at least one pair must match).
    These headers need to be specified with the {\tt queue\_bind}
    method as arguments.
\end{itemize}


\subsection{Predefined Exchanges}
With the program {\tt amq\_shell} \index{amq\_shell} you 
can find the following exchanges, which represent 
the types shown before. 
\begin{itemize}
  \item amq.fanout \index{amq.fanout}
  \item amq.direct \index{amq.direct}
  \item amq.topic  \index{amq.topic}
  \item amq.headers \index{amq.headers}   
\end{itemize}

There are several other objects shown with the {\tt ls}
command.


\section{Collaboration}

In this section a few structures are shown as an example how
the binding works.

\subsection{Structure of a program}
Most programs will have a structure shown below, 
calling the {\tt destroy} can be handled automatically
via the garbage collection in a safe way also.

\begin{itemize}
\item Create the WireAPI object,
\item create an AMQ\_Auth object with credentials
\item create an AMQ\_Connection object 
\item create an AMQ\_Session object
\item interact via the WireAPI
\item destroy the AMQ\_Session object
\item destroy the AMQ\_Connection object
\end{itemize}


\subsection{Creating a consumer}
The consumer receives data from the OpenAMQ server. 
There is also a basic\_get method but it is less 
efficient compared to a consumer.  It is not yet 
clear how it must be used. When binding a queue 
to an exchanges can be complex since you need to wait 
until the exchange exists. 
The predefined exchanges can be a solution for this
problem. It is also possible to test for the
existence of an exchange by using the {\tt passive} flag
when calling {\tt exchange\_declare}

\begin{itemize}
  \item create the AMQ\_Session object
  \item declare the queue, maybe just passive
  \item bind the queue to an exchange
  \item basic\_consume
  \item call wait
  \item call basic\_arrived to receive an AMQ\_ContentBasic
  \item repeat the last 2 steps until finished
  \item basic\_cancel
  \item delete the queue if needed
\end{itemize}


\subsection{Creating a publisher}
A publisher sends data to the OpenAMQ server. The server
will use the types in bindings to get the data flowing to
the right queues.

\begin{itemize}
  \item create the AMQ\_Session object
  \item declare the exchange, maybe only  passive
  \item create an AMQ\_ContentBasic object
  \item populate the content
  \item call basic\_publish to send it the the server
  \item delete the exchange if needed
\end{itemize}




\section{Classes}

A short list of all the methods of the object to show 
how to use them to interact with the server. The aim 
was to make it ruby style, not to translate all the names 
from the C functions. Since ruby has garbage collection, 
the most important classes support this by destroying them 
at them when they are remove by the garbage collector.

The WireAPI class is a bit special, since this is a
singleton, each instance will point to the same object.
When all instances are removed the internal structures
will also be removed.

\subsection{WireAPI}\index{WireAPI}

The namespace for this binding is all in this object. Before
using this binding make sure to have an object of the class
WireAPI since it existence controls the internal structures
of the WireAPI library. There are 2 methods to create debug
output that are for the entire WireAPI library.

\method{WireAPI.new}{number=0}{WireAPI}
       {
         This method will initialize the binding. The optional 
         argument defines the tracing level as a bitmap. Value
         4 controls the {\tt debug\_connection} and value 8 
         controls the {\tt debug\_session} these need to be ORed.      
         This parameter is only used when not 0 so to disable
         debugging use a value without bit values  4 and 8,
         for example '16'. These values were selected to not
         interfere with the values for the parameter trace
         of {\tt connection\_new} so tracing could be configured
         with a singele number.
         Try to have only one instance of this object, since all 
         ruby instances will point to the same singleton.         

         The tracing level can be changed at any time with the 
         {\tt debug\_connection=} and the {\tt debug\_session}
         methods.
       }


\method{auth\_plain}{username, password}{AMQ\_Auth}
       {
         creates an auth object to be used with the 
         client\_connection\_new method as the third argument. 
         Object of this type are destroyed by the garbage collector,
         it only exists because of the WireAPI library, it could
         have been implemented with 2 values in the method
         to create a connection.
       }


\method{connection\_new}
       {host, vhost, auth\_data, client\_name, trace, timeout}
       {AMQ\_Connection}
       {
         Creates a new connection to the server. On this connection
         the programmer needs to create an AMQ\_Session.
         An alternative style is to use a codeblock that receives the
         AMQ\_Connection as an argument at the end of this block the
         connection will be destroyed.
         \begin{itemize}
           \item host (string) - Specifies a server name 
             or IP address, optionally ending in ':' plus 
             a port number.
           \item vhost (string) - the vhost name can be nil for the default '/'
           \item client\_name (string) - The instance argument 
             sets the client instance name, which can be used 
             to identify a specific client in the management 
             console or server log.
           \item trace (number) - sets the trace level for WireAPI library,
             the lowest 2 bit will be used (0-3)
           \item timeout(number) - This argument governs 
             all synchronous exchanges with the server - if the 
             server does not respond within this time, the connection 
             treats it as a fatal error. A timeout of zero means 
             "infinite". A good value for fast networks is five 
             to ten seconds; for a slower network, a value of 30 
             seconds or more is reasonable. The unit of this value
             is in milleseconds.
         \end{itemize}
       }

\method{debug\_connection=}
       {flag}{self}
       {
         enable or disable tracing of the connection to stderr 
       }


\method{debug\_session=}
       {flag}{self}
       {
         enable or disable tracing of the session to stderr
       }


\subsection{AMQ\_Connection}\index{WireAPI::AMQ\_Connection}

A connection is a TCP/IP connection from a client 
to an OpenAMQ server. AMQP is a multi-channel protocol, 
meaning that one network connection can carry an arbitrary 
number \footnote{OpenAMQ currently supports exactly ONE 
session per connection.} of parallel, independent virtual 
connections, which AMQP calls "channels". In WireAPI these 
are called "sessions" for compatibility with other 
middleware APIs. 

\method{alive}{}{boolean}
       {returns false when connection has had an error}
       
\method{client\_session\_new}{}{AMQ\_session} 
       {Opens a session, in OpenAMQ there is only one session for
        a connection.
        An alternative style is to use a codeblock that receives the
        AMQ\_Session as an argument at the end of this block the
        session will be destroyed.
       }

\method{destroy}{}{nil} 
       {Destroys the connection, garbage collection will destroy the object
       also, but calling this method allows the programmer to trigger it,
       after this the internal object will no longer be usable.}


\method{error\_text}{}{string} 
       {error string reported by the API}

\method{reply\_code}{}{number} 
       {error value reported by server}

\method{reply\_text}{}{string} 
       {error string reported by the server}

\method{server\_copyright}{}{string} 
       {copyright notice reported by server}

\method{server\_platform}{}{string} 
       {operating system platform reported by server}

\method{server\_product}{}{string} 
       {product name reported by server}

\method{server\_version}{}{string} 
       {product version reported by server}

\method{silent}{}{boolean} 
       {suppresses error reporting when true}

\method{silent=}{boolean}{self} 
       {suppresses error reporting when set to true}



\subsection{AMQ\_Session}\index{WireAPI::AMQ\_Session}

A session corresponds to an AMQP channel, and is 
a virtual connection to an AMQP server. You must 
at least create one session in order to talk with 
an AMQP server. While AMQP offers multiplexing in 
theory, OpenAMQ no longer implements this, mainly 
because there are no proven performance advantages. 
So you should create a single session per connection, 
no more. Future versions of WireAPI may merge the 
session and connection classes.


When an AMQ\_Error\index{WireAPI::AMQ\_Error} is thrown, 
in most cases the error\_text will be used as message. 
Unless the message was not from the WireAPI library but 
from the binding. These strings don't offer a number code 
at this moment for easier handling. 
On the other hand the AMQ\_ServerError
\index{WireAPI::AMQ\_ServerError}  contains a reply\_code method
which makes it easier for a program to be aware of 
the cause of the exception.

\method{active}{}{boolean}
       {??}

\method{alive}{}{boolean}
       {returns false when connection has had an error.
       If the connection is broken this implies the session}

\method{basic\_ack}{delivery\_tag, multiple}{boolean}
       {
         \begin{itemize}
           \item delivery\_tag (AMQ\_int64) - can be retrieved with 
             the {\tt delivery\_tag} method
           \item multiple (boolean) -
         \end{itemize}
       }

\method{basic\_arrived}{}{AMQ\_Content}
       {
       returns nill when no content is available.
       returns the content, that is available after a call to {\tt wait}
       }


\method{basic\_cancel}{consumer\_tag}{self}
       {
       The consumer\_tag is a string the was used for the 
       basic\_consume. It ends the consumer
       }


\method{basic\_cancel\_nowait}{consumer\_tag}{self}
       {
       similar to basic\_cancel but doesn't wait until
       it is finished
       }


\method{basic\_consume}
       {queue\_name, consumer\_tag, no\_local, no\_ack, exclusive, arguments=nil}
       {self}
       {
       \begin{itemize}
         \item queue\_name (string) - the queue to receive data from
         \item consumer\_tag (string) - a possible empty string to identify
           this consumer
         \item no\_local (boolean) - messages from this session will 
           not be received when true
         \item no\_ack (boolean) -
         \item exclusive (boolean) -
         \item arguments (AMQ\_FieldList) - optional. A hash will be 
           converted to an AMQ\_FieldList when used.
       \end{itemize}
       after basic\_consume use wait and basic\_arrived
       to receive the data until calling basic\_cancel.
       }
∑  

\method{basic\_consume\_nowait}{queue\_name, consumer\_tag, no\_local, no\_ack,exclusive}
       {self}
       {
       similar to basic\_consume but doesn't wait until
       it is finished
       }


\method{basic\_get}{queue\_name, no\_ack}{self}
       {
         a way to receive without a consumer, the performance will 
         be very low because of synchronization. 
       }


\method{basic\_publish}
       {basic\_content, exchange\_name, routing\_key, mandatory, immediate, arguments=nil}
       {self}
       {
         \begin{itemize}
           \item basic\_content (WireAPI::AMQ\_ContentBasic) - the data to send
           \item exchange\_name (string) - name of an exchange on the server
           \item routing\_key (string) - a key used for routing in a dot 
             separate form, the type of the exchange defines how this 
             key is used
           \item mandatory (boolean) -  message must be routable
           \item immediate (boolean) -  message must be deliverable.
         \end{itemize}
       }

\method{basic\_reject}
       {delivery\_tag, requeue}
       {boolean}
       {
         \begin{itemize}
         \item delivery\_tag (AMQ\_int64) - can be retrieved with 
           the {\tt delivery\_tag} method
         \item requeue (boolean) -
         \end{itemize}
       }


\method{consumer\_count}{}{number}
       {
         returns the number of consumers
       }

\method{consumer\_tag}{}{number}
       {
         server generated consumer tag, after creating a consumer with
         {\tt basic\_consume}, when a consumer\_tag was specified in
         the {\tt basic\_consume} it will be returned here. Otherwise
         the server will create a value which will needed to cancel 
         the consumer.
       }

\method{delivery\_tag}{}{AMQ\_int64}
       {
         the delivery tag which can be used for basic\_ack and 
         basic\_reject. Even in the logging it shows 0 at this 
         moment. Which means the binding will also get always the
         same 0 value.
       }


\method{destroy}{}{nil} 
       {Destroys the session, garbage collection will destroy the object
       also, but calling this method allows the programmer to trigger it,
       after this the internal object will no longer be usable.}

\method{error\_text}{}{string} 
       {error string reported by the API}


\method{exchange}{}{string} 
       {exchange name from last method}

\method{exchange\_declare}
       {exchange\_name, type, passive, durable, auto\_delete, internal, arguments=nil}
       {self} 
       {
         create a new exchange. There are predefined exchanges called 
         AMQ.* but in case there is need for another one they 
         can be created.
         \begin{itemize}
         \item exchange\_name (string) - The exchange name is a 
           client-selected string that identifies the exchange 
           for publish methods. Exchange names may consist of 
           any mixture of digits, letters, and underscores. 
           Exchange names are scoped by the virtual host.
         \item type (string) - Each exchange belongs to one of 
            a set of exchange types implemented by the server. 
            The exchange types define the functionality of the 
            exchange - i.e. how messages are routed through it. 
            It is not valid or meaningful to attempt to change 
            the type of an existing exchange. OpenAMQ supports 
            "fanout", "direct", "topic", and "header" exchanges.
         \item passive (boolean) - do not create exchange. If set, 
            the server will not create the exchange. The client 
            can use this to check whether an exchange exists without 
            modifying the server state.
         \item durable (boolean) - not supported by OpenAMQ
         \item auto\_delete (boolean) - If set, the exchange is 
            deleted when all queues have finished using it.
         \item internal (boolean) - If set, the exchange may not 
           be used directly by publishers, but only when bound to 
           other exchanges. Internal exchanges are used to construct 
           wiring that is not visible to applications. 
         \item arguments (AMQ\_FieldList) - internal configdata to
           which is probably server specific. A hash will be 
           coverted to an AMQ\_FieldList when needed.
         \end{itemize}
       }

\method{exchange\_declare\_nowait}
       {exchange\_name, type, passive, durable, auto\_delete, internal, arguments=nil}
         {self} 
       {
       same as exchange\_declare but don't wait to complete.
       }


\method{exchange\_delete}
       {exchange\_name}{self} 
       {
       deletes an exchange. An exchange can be predefined or have
       an auto\_delete flag enabled, in both cases this method
       will not be needed.
       }


\method{exchange\_delete\_nowait}
       {exchange\_name}{self} 
       {
       same as exchange\_delete, but don't wait to complete.
       }

\method{message\_count}
       {}{number} 
       {
         number of messages in queue
% which queue , consumer ??
       }

\method{queue}
       {}{string} 
       {
       queue name assigned by server, used for {\tt queue\_declare} with
       a nil as queue\_name argument. This value will be needed to
       be able to bind the created queue or delete it again. When
       a name was specified for a {\tt queue\_declare} this name
       will be returned.
       }


\method{queue\_bind}
       {queue\_name, exchange\_name, routing\_key, arguments=nil}{self} 
       {
         Bind a queue to an exchange. Until a queue is bound it 
         will not receive any messages. In a classic messaging model, 
         store-and-forward queues are bound to a dest exchange 
         and subscription queues are bound to a dest\_wild exchange.         
         \begin{itemize}
           \item queue\_name (string) - Specifies the name of the queue 
             to bind. If the queue name is empty (nil), refers to the 
             current queue for the channel, which is the last declared 
             queue. Queue names may consist of any mixture of digits, 
             letters, and underscores.
           \item exchange\_name (string) - The exchange name 
             is a client-selected string that identifies the exchange 
             for publish methods. Exchange names may consist of any 
             mixture of digits, letters, and underscores. Exchange 
             names are scoped by the virtual host.
           \item routing\_key (string) - Specifies the routing key 
             for the binding. The routing key is used for routing 
             messages depending on the exchange configuration. Not 
             all exchanges use a routing key - refer to the specific 
             exchange documentation. If the routing key is empty and 
             the queue name is empty, the routing key will be the 
             current queue for the channel, which is the last 
             declared queue.
           \item arguments (AMQ\_FieldList) - can be used for exchanges
             with type 'headers'. A hash will be coverted to an
             AMQ\_FieldList when used.
         \end{itemize}
       }


\method{queue\_bind\_nowait}
       {queue\_name, exchange\_name, routing\_key, arguments=nil}{self} 
       {
         same as queue\_bind, but don't wait to complete.
       }


\method{queue\_declare}
       {queue\_name,passive,durable,exclusive,auto\_delete,arguments=nil}{self} 
       {
         \begin{itemize}
           \item queue\_name (string) - Queue names may consist 
             of any mixture of digits, letters, and underscores. 
             May be specified, or may be nil. If the 
             queue name is nil, the server creates and names 
             a queue and returns this. You can access the last 
             created queue from with the method queue; if you want 
             to create many queues, copy the returned name somewhere safe.
           \item passive (boolean) - do not create queue. If set, 
             the server will not create the queue. The client can 
             use this to check whether a queue exists without 
             modifying the server state.
           \item durable (boolean) - A durable message queue is 
             like a configured object: when the server (re)starts, 
             the message queue is present and active. 
             Temporary queues are destroyed when the server 
             shuts down. OpenAMQ doesn't support this.
           \item exclusive(boolean) - request an exclusive queue. 
             Exclusive queues may only be consumed from by the 
             current connection. Setting the 'exclusive' flag 
             always implies 'auto\_delete'.
           \item auto\_delete (boolean) - If set, the queue is 
             deleted when all consumers have finished using it. 
             Last consumer can be cancelled either explicitly 
             or because its channel is closed. If there was no 
             consumer ever on the queue, it won't be deleted.
           \item arguments (AMQ\_FieldList) - optional.  A hash 
             argument will be converted when needed
         \end{itemize}
       }

\method{queue\_declare\_nowait}
       {queue\_name, passive, durable, exclusive, auto\_delete, arguments=nil}
       {self} 
       {
         same as queue\_declare , but don't wait to complete.       
       }


\method{queue\_delete}
       {queue\_name}{self} 
       {
         When a queue is deleted any pending messages are sent 
         to a dead-letter queue if this is defined in the server 
         configuration, and all consumers on the queue are cancelled.
         \begin{itemize}
           \item queue\_name (string) - Specifies the name of 
             the queue to delete. If the queue name is nil, 
             refers to the current queue for the channel, 
             which is the last declared queue. Queue names may 
             consist of any mixture of digits, letters, and underscores.
         \end{itemize}         
       }

\method{queue\_delete\_nowait}
       {queue\_name}{self} 
       {
         same as queue\_delete , but don't wait to complete.       
       }


\method{queue\_purge}
       {queue\_name}{self} 
       {
         This method removes all messages from a queue. It 
         does not cancel consumers. Purged messages are 
         deleted without any formal "undo" mechanism.
         \begin{itemize}
         \item queue\_name (string) - Specifies the name of the queue 
           to purge. If the queue name is empty, refers to the current 
           queue for the channel, which is the last declared queue. 
           Queue names may consist of any mixture of digits, letters, 
           and underscores.
         \end{itemize}

       }

\method{queue\_purge\_nowait}
       {queue\_name}{self} 
       {
         same as queue\_purge , but don't wait to complete.       
       }


\method{queue\_unbind}
       {queue\_name, echange\_name, routing\_key, arguments=nil}{self} 
       {
         Unbind a queue from an exchange.
         \begin{itemize}
           \item queue\_name (string) - Specifies the name 
             of the queue to unbind. If the queue name is 
             empty, refers to the current queue for the 
             channel, which is the last declared queue. 
             Queue names may consist of any mixture of digits, 
             letters, and underscores.
           \item exchange\_name (string) - The exchange 
             name is a client-selected string that identifies 
             the exchange for publish methods. Exchange names 
             may consist of any mixture of digits, letters, 
             and underscores. Exchange names are scoped by 
             the virtual host.
           \item routing\_key (string) - Specifies the routing key 
             of the binding to unbind.
           \item arguments (AMQ\_FieldList) - optional, depends on the
             type of exchange. A hash argument wile be converted when needed.
         \end{itemize}
       }

\method{queue\_unbind}
       {queue\_name, echange\_name, routing\_key, arguments=nil}{self} 
       {
         same as queue\_unbind , but don't wait to complete.       
       }


\method{redelivered}
       {}{boolean} 
       {
         message is being redelivered
       }

\method{reply\_code}{}{number} 
       {
         error value reported by server. When this value is not
         0 most methods will throw a AMQ\_ServerError.
       }


\method{reply\_text}{}{string} 
       {
         error string reported by the server. This will be used
         when AMQ\_ServerError is thrown as message.         
       }


\method{routing\_key}{}{string} 
       {
        the routing key from the last message.
       }

\method{wait}{number=0}{self} 
       {
         Waits for content to arrive from the server. You must call 
         this method in order to get content. Returns either because 
         content arrived or because the timeout (in milliseconds) 
         expired. This timeout value is optional, when no timeout is 
         specified (also for values 0 and nil) there is no timeout at all.
         The wait method should not be replaced by a OS sleep of any 
         type since {\tt wait} is aware of the session being alive and
         stops also when the {\tt basic\_arrive} method will return
         a result. Als the internals of WireAPI will not be blocked
         by it.
       }


\subsection{AMQ\_ContentBasic}\index{WireAPI::AMQ\_ContentBasic}

AMQP uses the term "content" to mean an application 
message (the term "message" means different things 
at the application, protocol, and internal technical 
levels, so is confusing).

These AMQ\_Contents can be routed to consumers in 
different ways, for each type the right exchange will
be chosen to public the content.

\begin{itemize}
  \item fanout - all bound subscribers receive it
  \item routing key - can handled in different ways (direct, topic)
  \item headers - the headers control which queue will receive it
\end{itemize}

Apart from these headers are some properties which are not
for routing but they can also contain various information.


\method{WireAPI::AMQ\_ContentBasic.new}{}{AMQ\_ContentBasic} 
       {
         creates a content object to be published via a session
       }


\method{app\_id}{}{string} 
       {
         Get creating application id.
       }


\method{app\_id=}{string}{self} 
       {
         sets creating application id.
       }


\method{body}{number=nil}{string} 
       {
         Get the body of the message as a string the maximal
         length can be specified optionally. Nil can also
         be used as maximal value. In most cases the default
         nil will be the best choice.
       }

\method{body=}{string}{self}
       {
         sets the body of the message 
       }

\method{body\_size}{}{number}
       {
         Get the body size of the content
       }

\method{content\_encoding}{}{string}
       {
         Get MIME content encoding.
       }

\method{content\_encoding=}{string}{}
       {
         Set MIME content encoding. 
       }

\method{content\_type}{}{string}
       {
         Get MIME content type.
       }

\method{content\_type=}{string}{}
       {
         Set MIME content type. 
       }


\method{correlation\_id}{}{string}
       {
         Get the correlation id, this can refer
         to the message id that triggered this
         message. Like the answer to a question
         with a message\_id.
       }

\method{correlation\_id=}{string}{}
       {
         Set the correlation id probably, this can be
         the message id from the received message that
         triggered sending a message.
       }

\method{delivery\_mode\_persistent}{}{boolean}
       {
         in the WireAPI library delivery mode can be non-persistent(1) 
         or persistent(2) this method returns this as a boolean
       }


\method{delivery\_mode\_persistent=}{boolean}{self}
       {
         in the WireAPI library delivery mode can be non-persistent(1) 
         or persistent(2) this method sets the value based on its
         argument.
       }

\method{destroy}{}{nil}
       {
         The same as {\tt unlink} method but this alias
         is similar to other objects. From the Ruby side the
         effect will be the same as the other {\tt destroy}
         methods.
       }

\method{exchange}{}{string}
       {
         Get the exchange\_name to which the content was published.
       }


\method{expiration}{}{string}
       {
         'Message expiration specification'
       }


\method{expiration=}{string}{self}
       {
         The format of this 'Message expiration specification' might
         must be uniform, but until today no specification has been
         found.
       }


\method{headers}{}{AMQ\_FieldList}
       {
         Get the headers of a message. When you are certain that
         all keys are unique this can be converted to a hash.
       }

\method{headers=}{AMQ\_FieldList}{}
       {
         Sets the headers of a message. The argument can also
         be a ruby Hash, in this case a new AMQ\_FieldList will 
         be created and the add\_hash method will be used to 
         populate the Fieldlist with the keys and values of 
         the hashtable.
       }


\method{message\_id}{}{string}
       {
         Get application message identifier of the object.
       }

\method{message\_id=}{string}{}
       {
         Sets the message identifier of the object.
       }

\method{priority}{}{number}
       {
         Get the message priority.
       }

\method{priority=}{number}{}
       {
         Set the message priority from in the range from 0 to 9.
       }

\method{producer\_id}{}{string}
       {
         Get the producer id of the publishing connection.
       }

\method{reply\_to}{}{string}
       {
         Get the destination to reply to.
       }

\method{reply\_to=}{string}{}
       {
         Set the destination to reply to.
       }

\method{routing\_key}{}{string}
       {
         The original routing\_key specified by the publisher.
       }

\method{type}{}{string}
       {
         Get the message type name. This will be in most cases
         a high level type, different from the content\_type.
       }

\method{type=}{string}{}
       {
         Set the message type name.
       }

\method{unlink}{}{nil}
       {
         Unlinks the internal object, this method will be
         called also by the garbage collector. After calling
         this method the object is no longer usable. The 
         internal object might still exist in another thread
       }

\method{user\_id}{}{string}
       {
         Get creating user id.
       }

\method{user\_id=}{string}{}
       {
         Set creating user id.
       }




\subsection{AMQ\_FieldList}\index{WireAPI::AMQ\_Connection}

The fieldlist tries to be similar to a hash, it can be 
converted to a hash in most cases. It can be also populated 
from a hash with the add\_hash method. The {\tt header=} method 
can also handle a hash instead of AMQ\_FieldList instances. 
The main difference with a hash is that the same key can 
appear multiple times. When this feature makes sense, using 
a hash will not be an option. The {\tt each\_pair} method will give
access to the keys and values when a key occurs multiple times.

The stored keys will be always a string, when they need to be
converted by calling their {\tt to\_s} method, this will be
called automatically.

\method{WireAPI::AMQ\_FieldList.new}{[AMQ\_FieldList|Hash]}{AMQ\_FieldList} 
       {
         creates a FieldList to be used with {\tt AMQ\_ContentBasic\#headers=}
         or as argument list for several method of {\tt AMQ\_Session}.
         This constructor can accept a hash as argument or work like a
         copy constructor.
       }


\method{add\_hash}{hash}{self} 
       {
         adds the keys from a hash with the value to the fieldlist. This is
         like using {\tt put} for all keys in the hash argument.
       }


\method{destroy}{}{nil} 
       {Destroys the fieldlist, garbage collection will destroy the object
       also, but calling this method allows the programmer to trigger it,
       after this the internal object will no longer be usable.}


\method{each\_pair}{}{self} 
       {
         iterates over all keys (also multiple ones) with a codeblock that
         receives 2 arguments, the key and the value. This is intented
         when not all keys are unique.
       }


\method{get}{string}{object} 
       {
         retrieves the value of a key or nil, in case of not unique keys only
         the first value will be returned.
       }


\method{keys}{}{array of strings}
       {
         returns an array of the keys similar to a hash. but multiple keys will
         appear more times, so this is a way to detect them. 
       }


\method{keys\_unique?}{}{boolean}
       {
         Returns true when all keys are unique. In this case converting
         to a hash will be no problem. When the keys are not unique
         converting to hash will loose the double keys.
       }


\method{put}{string, object}{self} 
       {
         add a key with a value to the list, the object can be currently
         5 types (string, Fixnum, Time, Float, or nil) The float numbers
         will be represented as fixed points, which implies a limited
         range. When the key is not a string, the method {\tt to\_s} will
         be called on the key. When key already existed before, there
         will be two instances of the key.
       }


\method{to\_hash}{}{hash} 
       {
         Transforms to a hash. When the keys are not unique the last
         value of the multiple keys will be used. Before calling this
         method a test with {\tt keys\_unique?} is a way to know
         it can be converted to a hash.
       }


\method{[]}{string}{object} 
       {
         an alias for the get method, to look more like a hash
       }


\method{[]=}{string, object} {self}
       {
         an alias for the put method, to look more like a hash.
         Be aware it add a key, it doesn't replace an existing key.
       }

\subsection{AMQ\_ServerError}\index{WireAPI::AMQ\_ServerError}
This is a subclass of the AMQ\_Error.\index{WireAPI::AMQ\_Error}

When the reply\_code of the server is not 0 after calling
a session method an object of this type will be thrown. The 
{\tt message} is the {\tt reply\_text} of the session and 
the {\tt reply\_code} can be retrieved with a method of 
the same name.


\method{message}{} {string}
       {
         return the reply\_text of the session since this
         will be a human readable explanation.
       }

\method{reply\_code}{} {number}
       {
         returns the code of the error, which will be easier to
         process than the human readable message. This will be
         similar to the {\tt AMQ\_Session\#reply\_code}
       }


\printindex

\end{document}

